% !TeX program = lualatex
% HTML (lwarp) variant of pyscfBook.tex. Sources of the HTML pipeline live in html/; html/build-html.ps1 copies this file
% to the project root for the build and removes it (and everything lwarp makes) afterwards.
\documentclass{PySCFBookHTML}
\addbibresource{../Bibliography/QuantumChemistry.bib}
\addbibresource[label=soft]{Articles.bib}

\title{Квантово-механічні методи обчислення}
\author{С.~М.~Пономаренко}
\begin{document}
\maketitle
\tableofcontents
\ForceHTMLPage
\htmlinput{intro}

\newcommand{\ChaptersI}{WelcomeWithPyscf,BasePyscf}
\newcommand{\ChaptersII}{HFMethod,HFMethodDiatomic,PostHFMethods,DFTMethod}
\newcommand{\ChaptersIII}{ResponseTheory,GeometryOptimization,Oscilations,ExcitedStates}

\part{Вступ до PySCF: основи та перші кроки}
\epigraph{«Теорія без практики --- порожня, практика без теорії — сліпа.»}{--- Іммануїл Кант, адаптовано}
\multiinclude{\ChaptersI}[]

\part{Ab Initio методи: від Хартрі-Фока до пост-HF і DFT}
\epigraph{«Рівняння --- це мова, якою природа говорить із нами.»}{--- Поль Дірак}
\multiinclude{\ChaptersII}[]

\part{Властивості молекул}
\epigraph{«Ми змушені визнавати, що кожне спостереження неминуче порушує явище, але водночас саме спостереження й дає нам знання про нього.»}{--- Нільс Бор}
\multiinclude{\ChaptersIII}[]

\appendix
\htmlinput{Additions/Abbr}
\htmlinput{Additions/mo-representation}
\htmlinput{Additions/DIIS}
\htmlinput{Additions/Convergence}

\printbibheading[heading=bibintoc]
\addtocategoryNoCi{pyscf}{sun2018pyscf, sun2020pyscf, sun2017pyscf}
\printbibliography[category=pyscf,title={Ключові статті по \PySCF},heading=subheading,resetnumbers=true]
\addtocategoryNoCi{Main}{YatsimirsijQChem, Sleta, Sun2025PythonQC, Jensen, Levine, Szabo}
\printbibliography[category=Main,title={Основна література},heading=subheading,resetnumbers=true]
\addtocategoryNoCi{Additional}{BSE, SQM, EGMI1, EGMI2, EGMI3}
\printbibliography[category=Additional,title={Додаткові посилання},heading=subheading,resetnumbers=true]
\end{document}
